\chapter{VR-Audit}
\label{chap:audit}
\uses{chap:apparatus, chap:sets}

\section{Position}
\label{sec:audit_position}

VR-Audit is the first Application in the VR Cycle (Layer~III).  It applies
the two-register apparatus from Chapter~\ref{chap:forms} and the Mode~B
schema from Chapter~\ref{chap:apparatus} to classical mathematics.  Audit One
establishes an operational Hahn-Banach theorem for Hilbert spaces via the
Riesz representation route (Path~B).

\textbf{Wrapping principle.} No new type of ``computable Hilbert space'' is
introduced parallel to mathlib's \texttt{Real} or \texttt{InnerProductSpace}.
Instead, a predicate \texttt{IsComputableReal : ℝ → Prop} selects the
computable reals as a sub-collection of mathlib's classical~$\mathbb{R}$.
Similarly, \texttt{OperationalHilbertSpace E} layers three computability
fields on top of mathlib's \texttt{InnerProductSpace ℝ E}.  This is the
direct embodiment of the VR-Forms two-register apparatus in analysis:
formal register $=$ mathlib's classical structures; operational register $=$
the sub-collection satisfying computability predicates with witnesses.

\textbf{Non-chronological dependency.} The Apparatus preprint
(Chapter~\ref{chap:apparatus}) was published after VR-Audit, but
methodologically, VR-Audit is an instance of the Mode~B schema defined
there.  The blueprint shows this dependency explicitly; publication order
is not the ordering principle here.

\noindent\axiomprofile{[propext, Classical.choice, Quot.sound]} for all
17 public objects.  \texttt{Classical.choice} enters through mathlib's
$\mathbb{R}$ (Cauchy completeness) and Riesz representation.  This is
the expected and acceptable ceiling: operational content lives within
the classical structure, not outside it.

%──────────────────────────────────────────────────────────────────────
\section{Operational structures (Stages 1--4)}
\label{sec:audit_structures}

\begin{definition}[Computability predicate for real numbers]
\label{def:IsComputableReal}
\lean{VR.Audit.IsComputableReal}
\leanok
$\texttt{IsComputableReal}(x) := \exists\, (\texttt{alg} : \mathbb{N}
\to \mathbb{Q})\, (\texttt{mod} : \mathbb{N} \to \mathbb{N}),\;
\forall n\, k,\; \texttt{mod}(n) \le k \Rightarrow
|(\texttt{alg}(k) : \mathbb{R}) - x| \le 1/2^n$.
Following Pour-El \& Richards~1989~§\,1.1 and Bishop-Bridges~1985~§\,2.2:
$x$ is computable if there exists a rational approximation sequence with
an explicit modulus of convergence.  Implemented as a \texttt{Prop}
(not a \texttt{Structure}): witnesses are accessible via \texttt{obtain}
when needed.
\end{definition}

\begin{definition}[Computability predicate for sequences]
\label{def:IsComputableSequence}
\lean{VR.Audit.IsComputableSequence}
\leanok
\uses{def:IsComputableReal}
$\texttt{IsComputableSequence}(s : \mathbb{N} \to \mathbb{R}) :=
\forall n,\; \texttt{IsComputableReal}(s(n))$: every term of the
sequence is a computable real.
\end{definition}

\begin{theorem}[Computability closure lemmas]
\label{thm:IsComputableReal_lemmas}
\lean{VR.Audit.IsComputableReal_rat, VR.Audit.IsComputableReal_zero,
      VR.Audit.IsComputableReal_one, VR.Audit.IsComputableReal_neg,
      VR.Audit.IsComputableReal_add, VR.Audit.IsComputableReal_sub}
\leanok
\uses{def:IsComputableReal}
Rational numbers are computable (\texttt{IsComputableReal\_rat});
$0$, $1$ are computable; the predicate is closed under negation,
addition, and subtraction.  These six lemmas at
\axiomprofile{[propext, Classical.choice, Quot.sound]}.
Note: \texttt{IsComputableReal\_mul} is absent from Stage~1 (requires
bounded-sequence argument handled at the Hilbert-space level); it is
not needed for the main theorem.
\end{theorem}

\begin{definition}[Operational Hilbert space]
\label{def:OperationalHilbertSpace}
\lean{VR.Audit.OperationalHilbertSpace}
\leanok
\uses{def:IsComputableReal}
\texttt{OperationalHilbertSpace~E} is a typeclass (not a structure)
requiring \texttt{[InnerProductSpace ℝ E]} and \texttt{[CompleteSpace E]}
as implicit prerequisites, with three additional fields:
\begin{enumerate}
\item \texttt{denseSeq : ℕ → E} --- an explicit dense sequence
  (separability witness);
\item \texttt{denseSeq\_dense : DenseRange denseSeq} --- its range is
  classically dense in $E$;
\item \texttt{inner\_computable : ∀ m n, IsComputableReal (@inner ℝ E \_
  (denseSeq m) (denseSeq n))} --- pairwise inner products of the dense
  sequence are computable reals.
\end{enumerate}
Three fields are the minimum sufficient for the main theorem
(Stage~5 via Riesz).  No completeness modulus is required: classical
\texttt{CompleteSpace E} from mathlib suffices.  Pointwise computability
of individual \texttt{denseSeq k} is not required.
\end{definition}

\begin{definition}[Operational located subspace]
\label{def:OperationalLocatedSubspace}
\lean{VR.Audit.OperationalLocatedSubspace}
\leanok
\uses{def:OperationalHilbertSpace}
$\texttt{OperationalLocatedSubspace E}$ is a structure extending
\texttt{ClosedSubmodule ℝ E} (mathlib) with three fields:
\begin{enumerate}
\item \texttt{denseSubSeq : ℕ → E} --- explicit dense sequence within~$M$;
\item \texttt{denseSubSeq\_dense\_in} --- density within $M$;
\item \texttt{dist\_computable : ∀ n, IsComputableReal (infDist
  (denseSeq n) M)} --- computable distance from ambient dense sequence
  points to~$M$ (locatedness in Bishop's sense).
\end{enumerate}
\texttt{HasOrthogonalProjection} is inherited from
\texttt{ClosedSubmodule} via a global mathlib instance; no extra field
is needed.  Locatedness is formulated only over \texttt{denseSeq}
points, not arbitrary \texttt{x : E}: see Observation~3.
\end{definition}

\begin{definition}[Operational normable functional]
\label{def:OperationalNormableFunctional}
\lean{VR.Audit.OperationalNormableFunctional}
\leanok
\uses{def:OperationalLocatedSubspace, def:IsComputableReal}
$\texttt{OperationalNormableFunctional~E~M}$ is a structure wrapping
a classical \texttt{M.toSubmodule →L[ℝ] ℝ} (continuous linear map)
with two operational fields:
\begin{enumerate}
\item \texttt{fn\_computable\_on\_dense : ∀ n, IsComputableReal
  (toFun (denseSubSeq n, \_))} --- values on $M$'s dense sequence
  are computable;
\item \texttt{norm\_computable : IsComputableReal toFun.opNorm} ---
  the operator norm is a computable real (\emph{normability}).
\end{enumerate}
Distinction \emph{bounded} vs \emph{normable} (Ishihara~1989): in
computable analysis, a functional can be bounded without its norm being
computable.  \texttt{norm\_computable} captures exactly normability,
which is required for the norm-preservation clause of the main theorem.
See Observation~5.
\end{definition}

\begin{theorem}[Computability extends to all of $M$]
\label{thm:fn_computable_everywhere}
\lean{VR.Audit.fn_computable_everywhere}
\leanok
\uses{def:OperationalNormableFunctional, def:OperationalLocatedSubspace}
For any $\texttt{f : OperationalNormableFunctional~E~M}$ and
$x \in M.\texttt{toSubmodule}$, \texttt{IsComputableReal}$(f(x))$.
Proof: density gives a sequence $y_k \to x$; continuity
(\texttt{f.le\_opNorm}) bounds $|f(x) - f(y_k)|$;
\texttt{fn\_computable\_on\_dense} provides rational approximations
for $f(y_k)$; triangle inequality combines them.  Witnesses via
\texttt{Classical.choose}.
\end{theorem}

%──────────────────────────────────────────────────────────────────────
\section{Main theorem}
\label{sec:audit_main}

\begin{theorem}[Operational Hahn-Banach for Hilbert spaces]
\label{thm:HahnBanachOperational_Hilbert}
\lean{VR.Audit.HahnBanachOperational_Hilbert}
\leanok
\uses{def:OperationalNormableFunctional, def:OperationalLocatedSubspace,
      def:OperationalHilbertSpace, thm:fn_computable_everywhere}
Given: $E$ with \texttt{OperationalHilbertSpace E};
$M : \texttt{OperationalLocatedSubspace E}$;
$f : \texttt{OperationalNormableFunctional E M}$.

Produces: $g : E \to_L[\mathbb{R}] \mathbb{R}$ (classical CLM) such that:
(1)~$\forall n,\; \texttt{IsComputableReal}(g(\texttt{denseSeq n}))$
--- $g$ is computable on ambient dense sequence;
(2)~$\texttt{IsComputableReal}(g.\texttt{opNorm})$ --- $g$ is normable;
(3)~$\forall x \in M,\; g(x) = f(x)$ --- $g$ extends $f$;
(4)~$g.\texttt{opNorm} = f.\texttt{opNorm}$ --- norms agree.

Proof via Riesz representation (Path~B).  Six steps: (i)~Riesz vector
$\xi \in M$ via \texttt{InnerProductSpace.toDual}; (ii)~extension
$g := \texttt{innerSL}\,\mathbb{R}\,(\xi : E)$; (iii)~extension property;
(iv)~operationality of $g$ via factorisation $g(x) = f(P_M(x))$;
(v)~norm equality by Cauchy-Schwarz and inner product evaluation;
(vi)~computability of $g.\texttt{opNorm}$ from $f.\texttt{norm\_computable}$.
\noindent\axiomprofile{[propext, Classical.choice, Quot.sound]}.
\end{theorem}

%──────────────────────────────────────────────────────────────────────
\section{The transit pattern in action}
\label{sec:audit_transit}

\subsection*{Step~1: Riesz invoked as black box}

\begin{remark}[Riesz as formal-register tool]
\label{rem:riesz_black_box}
\lean{VR.Audit.HahnBanachOperational_Hilbert}
\leanok
\uses{thm:HahnBanachOperational_Hilbert}
\texttt{InnerProductSpace.toDual} (Riesz representation, mathlib) is
invoked to produce a vector $\xi \in M$ satisfying $f(v) = \langle\xi,
v\rangle_M$ for all $v \in M$.  $\xi$ has classical ontological status:
it exists by Riesz but its value is not explicitly computed.  This is the
formal-register step of the transit: a classical existence result is invoked
as a black box.  It uses \texttt{Classical.choice} through
\texttt{CompleteSpace} and \texttt{toDual}.
\end{remark}

\subsection*{Step~2: Explicit extension}

\begin{remark}[Concrete extension]
\label{rem:explicit_extension}
\lean{VR.Audit.HahnBanachOperational_Hilbert}
\leanok
$g := \texttt{innerSL}\,\mathbb{R}\,(\xi : E) : E \to_L[\mathbb{R}]
\mathbb{R}$ is defined concretely from $\xi$ without additional classical
machinery.  $g(w) = \langle\xi, w\rangle_E$ for all $w : E$.  Linearity
and continuity are automatic from \texttt{innerSL}.
\end{remark}

\subsection*{Step~3: Operationality via orthogonal projection}

\begin{remark}[Factorisation as operational heart]
\label{rem:orth_proj_transit}
\lean{VR.Audit.HahnBanachOperational_Hilbert}
\leanok
\uses{thm:fn_computable_everywhere, def:OperationalLocatedSubspace}
Key factorisation: $g(\texttt{denseSeq n}) = \langle\xi, \texttt{denseSeq
n}\rangle_E = \langle\xi, P_M(\texttt{denseSeq n})\rangle_M
= f(P_M(\texttt{denseSeq n}))$, where $P_M$ is orthogonal projection
(available from \texttt{HasOrthogonalProjection} inherited via
\texttt{ClosedSubmodule}).  Locatedness of $M$ (\texttt{dist\_computable})
makes the projection operationally accessible.  \texttt{fn\_computable\_everywhere}
gives computability of $f$ on all of $M$, including $P_M(\texttt{denseSeq n})$.
The transit closes: the formal object $\xi$ is used as an instrument to
express the operational result $g(\texttt{denseSeq n})$ as $f$ applied
to a computable input.
\end{remark}

\subsection*{Step~4: Norm equality}

\begin{remark}[Norm preservation]
\label{rem:norm_equality}
\lean{VR.Audit.HahnBanachOperational_Hilbert}
\leanok
$g.\texttt{opNorm} = \|\xi\|_E$ via \texttt{innerSL\_apply\_norm}.
$f.\texttt{opNorm} = \|\xi\|_E$ by \texttt{le\_antisymm}: upper bound
via Cauchy-Schwarz; lower bound by evaluating $f.\texttt{le\_opNorm}$
at $\xi$.  Transitively: $g.\texttt{opNorm} = f.\texttt{opNorm}$.
Computability of $g.\texttt{opNorm}$ follows from
$f.\texttt{norm\_computable}$ and the equality.
\end{remark}

%──────────────────────────────────────────────────────────────────────
\section{Why Path~B, not Path~A}
\label{sec:audit_path_b}

\begin{remark}[Specker obstruction and structural avoidance]
\label{rem:path_b}
The Specker obstruction (1949) applies to suprema of bounded monotone
computable sequences: such suprema can be non-computable.  Path~A ---
invoking classical Hahn-Banach directly --- produces an extension
\emph{abstractly}: only an existence result, with no computable witness
for $g(\texttt{denseSeq n})$.  This satisfies the formal register but
leaves the operational register empty.

Path~B (via Riesz) avoids the obstruction structurally:
(i)~$\xi$ is a \emph{single vector} produced by Riesz, not a limit of a
  monotone sequence;
(ii)~$g(\texttt{denseSeq n}) = f(P_M(\texttt{denseSeq n}))$ is a
  finite-step computation, not a sequence supremum;
(iii)~$g.\texttt{opNorm} = \|\xi\|_E$, a single real, not a supremum.

Locatedness (\texttt{dist\_computable}) is the key structural ingredient:
it makes orthogonal projection applicable at the operational level and
enables the factorisation of Step~3.  Without locatedness, the projection
$P_M(\texttt{denseSeq n})$ would remain in the formal register.
\end{remark}

%──────────────────────────────────────────────────────────────────────
\section{Non-vacuity instance}
\label{sec:audit_euclidean}

\begin{theorem}[$\mathrm{EuclideanSpace}\,\mathbb{R}\,(\mathrm{Fin}\,n)$ is an operational Hilbert space]
\label{thm:instEuclidean}
\lean{VR.Audit.instOperationalHilbertSpaceEuclidean}
\leanok
\uses{def:OperationalHilbertSpace, def:IsComputableReal}
For every $n : \mathbb{N}$, \texttt{instOperationalHilbertSpaceEuclidean}
provides \texttt{OperationalHilbertSpace (EuclideanSpace ℝ (Fin n))}.
Dense sequence: rational vectors $(\mathrm{Fin}\,n \to \mathbb{Q})$
enumerated via \texttt{Encodable.decode}, embedded via \texttt{WithLp.toLp 2}.
Density: three-step \texttt{DenseRange.comp} chain via
\texttt{Encodable.surjective\_decode\_getD},
\texttt{DenseRange.piMap}, and \texttt{PiLp.homeomorph}.
Inner product computability: \texttt{PiLp.inner\_apply} reduces $\langle\cdot,
\cdot\rangle$ to a finite sum of rational products; \texttt{IsComputableReal\_rat}
closes.

This is the canonical non-vacuity witness: for every $n$,
\texttt{HahnBanachOperational\_Hilbert} applies to
$\mathrm{EuclideanSpace}\,\mathbb{R}\,(\mathrm{Fin}\,n)$.
The infinite-dimensional case $\ell^2$ is deferred (requires
\texttt{Finsupp}-based enumeration and \texttt{tsum} convergence).
\end{theorem}

%──────────────────────────────────────────────────────────────────────
\section{Open programme}
\label{sec:audit_open}

\begin{remark}[Future audit candidates]
\label{rem:open_programme}
VR-Audit is an open programme.  Audit One (Hahn-Banach for Hilbert spaces
via Riesz) establishes the pattern.  Future candidates listed in the
preprint: Banach-Steinhaus (uniform boundedness), open mapping theorem,
closed graph theorem, spectral theorem for self-adjoint compact operators,
and Stone-Weierstrass.  Each would follow the Mode~B schema: classical
theorem invoked as formal-register black box; operational output extracted
via explicit witnesses and locatedness/computability predicates.
\end{remark}

%──────────────────────────────────────────────────────────────────────
\section{Position relative to neighbouring frameworks}
\label{sec:audit_comparison}

\begin{remark}[Comparison with Bishop-style constructive analysis]
\label{rem:compare_bishop}
Bishop-Bridges~1985 rebuilds analysis from scratch with constructive proofs.
VR-Audit wraps over mathlib's classical analysis rather than rebuilding it.
The operational layer is thinner (predicate restrictions, not new types)
and requires fewer proof lines, at the cost of accepting
\texttt{Classical.choice} in the formal register.
\end{remark}

\begin{remark}[Comparison with reverse mathematics]
\label{rem:compare_reverse_math}
Reverse mathematics locates each theorem in the subsystem hierarchy
(WKL$_0$, ACA$_0$, \ldots).  VR-Audit does not attempt this classification;
it identifies the \emph{operational content} that can be extracted from a
classical theorem at the \texttt{Classical.choice} ceiling, not the minimal
axioms needed to prove it.
\end{remark}

\begin{remark}[Comparison with computable analysis and Weihrauch reducibility]
\label{rem:compare_weihrauch}
Pour-El \& Richards~1989 and the Weihrauch framework study computability of
classical analysis results over abstract computation models.  VR-Audit works
within Lean~4's type-theoretic framework: computability is expressed as
\texttt{IsComputableReal} predicates rather than oracle Turing machines.
The metatheoretic equivalence of these notions is not formalised here.
\end{remark}

%──────────────────────────────────────────────────────────────────────
\section{Methodological observations}
\label{sec:audit_method}

\begin{remark}[Obs~1 --- No \texttt{Computable} annotation on approximation functions]
\label{rem:audit_obs1}
\lean{VR.Audit.IsComputableReal}
\leanok
\uses{def:IsComputableReal}
\texttt{Computable₂} for $\mathbb{Q}$ arithmetic operations is absent from
mathlib4 (verified by exhaustive search).  Lean's intrinsic totality of
\texttt{alg : ℕ → ℚ} and \texttt{mod : ℕ → ℕ} provides algorithmicity at
the type level; Turing-machine codings remain metatheoretic.  This is the
first concrete confirmation in VR-Audit of the architectural principle:
operationality is expressed through predicate restrictions, not
Lean's computational machinery.  (Parallel to VR-Numbers~§\,IV.1 boundary.)
\end{remark}

\begin{remark}[Obs~2 --- \texttt{inner\_computable} as wrapping principle in Hilbert setting]
\label{rem:audit_obs2}
\lean{VR.Audit.OperationalHilbertSpace}
\leanok
\uses{def:OperationalHilbertSpace}
The field \texttt{inner\_computable : ∀ m n, IsComputableReal
(@inner ℝ E \_ (denseSeq m) (denseSeq n))} selects, from the
classical values provided by \texttt{InnerProductSpace ℝ E}, precisely
those pairs for which the inner product is computable.  The operational
register (computable pairs) lives \emph{inside} the formal register
(full classical \texttt{InnerProductSpace}).  This generalises: any
classical structure with values in $\mathbb{R}$ can be wrapped operationally
by adding \texttt{IsComputableReal} predicates on selected outputs.
\end{remark}

\begin{remark}[Obs~3 --- Locatedness over dense sequence points only]
\label{rem:audit_obs3}
\lean{VR.Audit.OperationalLocatedSubspace}
\leanok
\uses{def:OperationalLocatedSubspace}
\texttt{dist\_computable} is formulated only over \texttt{denseSeq n}
(operational witnesses), not arbitrary $x : E$.  This is principled:
operational predicates apply to operational objects.  Arbitrary $x : E$
lives in the formal register.  Demanding computability on arbitrary
points would conflate the registers.  The formulation is sufficient
for Stage~5 because orthogonal projection is applied precisely to
\texttt{denseSeq} points.  This pattern --- operational predicates only
over operational witnesses --- generalises to all future VR-Audit work.
\end{remark}

\begin{remark}[Obs~4 --- Specker boundary structurally avoided]
\label{rem:audit_obs4}
\lean{VR.Audit.HahnBanachOperational_Hilbert}
\leanok
\uses{thm:HahnBanachOperational_Hilbert}
The Riesz vector $\xi$ is a single element, not a limit of a bounded
monotone computable sequence.  The Specker obstruction (1949) applies to
the latter; our construction explicitly avoids it.  Locatedness enables
orthogonal projection without a Specker-vulnerable supremum.  This is the
structural reason why Path~B (Riesz) is clean where a general Banach
Hahn-Banach extension would not be.
\end{remark}

\begin{remark}[Obs~5 --- Bounded $\ne$ normable]
\label{rem:audit_obs5}
\lean{VR.Audit.OperationalNormableFunctional}
\leanok
\uses{def:OperationalNormableFunctional}
In computable analysis (Ishihara~1989), ``bounded'' ($\exists C,\; |f(x)|
\le C\|x\|$) and ``normable'' (exact $\sup |f(x)|/\|x\|$ computable) are
distinct: a functional can be bounded without its norm being computable.
\texttt{norm\_computable} captures normability; this is required for the
norm-preservation clause $g.\texttt{opNorm} = f.\texttt{opNorm}$ in the
main theorem.  The distinction collapses classically; it is critical
constructively.
\end{remark}

\begin{remark}[Obs~6 --- Transit pattern has richer structure]
\label{rem:audit_obs6}
\lean{VR.Audit.HahnBanachOperational_Hilbert}
\leanok
\uses{thm:HahnBanachOperational_Hilbert}
The classical extended functional $g$ is built entirely from operational
data: $\xi$ is the Riesz vector of $f$; $g := \texttt{innerSL}\,
\mathbb{R}\,(\xi : E)$ is defined concretely.  The operational content
of $g$ at \texttt{denseSeq n} comes not from $g$'s definition directly
but from the factorisation $g(x) = f(P_M(x))$.  This is richer than
the standard Mode~B schema (``apply classical lemma, prove output
operational''): here the classical construction and the operational
extraction are interleaved through the projection factorisation.
\end{remark}

\begin{remark}[Obs~7 --- Norm instance synthesis gap]
\label{rem:audit_obs7}
\lean{VR.Audit.OperationalNormableFunctional, VR.Audit.HahnBanachOperational_Hilbert}
\leanok
The \texttt{Norm (M.toSubmodule →L[ℝ] ℝ)} instance does not synthesize
at declaration/goal time in this context (mathlib infrastructure gap).
Workaround: use \texttt{f.opNorm} (direct field access on
\texttt{ContinuousLinearMap}) throughout, and \texttt{change}/\texttt{show}
to pass to $\|\cdot\|$-based mathlib lemmas inside tactic proofs.  This
gap is structural (universe/instance unification) and expected in
wrapping-style formalisations with multiple unification paths.
\end{remark}

\begin{remark}[Obs~8 --- \texttt{⟪·,·⟫\_ℝ} notation fails in class and binder scopes]
\label{rem:audit_obs8}
\lean{VR.Audit.OperationalHilbertSpace}
\leanok
\uses{def:OperationalHilbertSpace}
The notation \texttt{⟪·, ·⟫\_ℝ} (from \texttt{open scoped
RealInnerProductSpace}) does not parse inside \texttt{class where} field
type declarations (subscript \texttt{\_ℝ} misinterpreted by the parser)
nor in $\exists$-binder scopes within tactic proofs.  Workaround in
class fields: use \texttt{@inner ℝ E \_} (explicit arguments).
Workaround in binder scopes: use \texttt{inner ℝ} (function form).
The notation works normally outside class blocks and binder scopes.
This is the second cross-namespace universe issue in VR-Audit (after
the \texttt{OSet.\{0\}} annotation in VR-Forms).
\end{remark}

\begin{remark}[Obs~9 --- \texttt{letI} vs \texttt{haveI} opacity in Riesz invocation]
\label{rem:audit_obs9}
\lean{VR.Audit.HahnBanachOperational_Hilbert}
\leanok
\uses{thm:HahnBanachOperational_Hilbert}
In the main theorem proof, \texttt{hK\_complete} (\texttt{CompleteSpace
K}) must be introduced via \texttt{letI} (transparent) rather than
\texttt{haveI} (opaque) so that tactic-level instance synthesis finds it.
Furthermore, \texttt{@InnerProductSpace.toDual} must be called with
explicit \texttt{hK\_complete} argument because instance synthesis fails
to match the unreduced form \texttt{↥M.toSubmodule} against
\texttt{↥↑M.toClosedSubmodule}.  These two obstacles are overcome without
new axioms; they reflect the instance-resolution sensitivity of
wrapping-style formalisations.
\end{remark}

\begin{remark}[Obs~10 --- Non-vacuity of operationality]
\label{rem:audit_obs10}
\lean{VR.Audit.instOperationalHilbertSpaceEuclidean}
\leanok
\uses{thm:instEuclidean}
Stage~6 confirms that the operational requirements of
\texttt{OperationalHilbertSpace} are achievable on a natural infinite
family: $\mathrm{EuclideanSpace}\,\mathbb{R}\,(\mathrm{Fin}\,n)$ for all
$n : \mathbb{N}$.  This is the standard VR-Audit pattern: main theorem
in full generality, then at least one explicit instance to show the
operationality is inhabited without exotic mathematics.
\end{remark}

%──────────────────────────────────────────────────────────────────────
\section{Axiom profile}
\label{sec:audit_axiom_profile}

\axiomprofile{[propext, Classical.choice, Quot.sound]} for all 17 public
objects.

\texttt{Classical.choice} enters through:
mathlib's $\mathbb{R}$ (Cauchy completeness, inherited by every object);
\texttt{InnerProductSpace.toDual} (Riesz, used in the main theorem);
\texttt{Classical.choose} in \texttt{fn\_computable\_everywhere}
(density approximation).

\begin{remark}[AC as formal-register tool]
\label{rem:AC_methodological}
\leanok
The use of \texttt{Classical.choice} in VR-Audit is methodologically
correct.  \texttt{Classical.choice} (AC) lives in the formal register:
it produces a classical existence result without a computable witness.
The operational register is then extracted via structural arguments
(projection factorisation, locatedness, normability) that do not
themselves require AC.  This is exactly the Mode~B pattern of
Chapter~\ref{chap:apparatus}: classical operations with operational
witnesses.  AC is a formal-register black box; the operational content
is the computable output extracted from it.
\end{remark}

%──────────────────────────────────────────────────────────────────────
\section{References}
\label{sec:audit_refs}

The preprint for this chapter is
\zenodo{10.5281/zenodo.20364111}{VR-Audit, v1.0.0} (24 May 2026).

The Lean~4 formalisation is archived at
\zenodo{10.5281/zenodo.20363739}{VR-Audit Lean, v1.0.0} (24 May 2026),
git tag \texttt{v1.4-vr-audit-hb-hilbert}.
